长直同轴导体的内外半径分别为 \(a,b\)（\(0<a<b\)），忽略端部效应；两导体间填充线性各向同性非均匀介质 \(\varepsilon(r)=\varepsilon_0(1+\alpha r/a)\)，其中 \(\alpha>0\)，内圆柱单位长度自由电荷为 \(+\lambda\)，外导体为 \(-\lambda\)。假定介质静电、无自由体电荷、真空磁导率。
(1) 求位移场、电场及两导体电势差，积分结果应明确依赖半径。
(2) 给出单位长度电容和场能，验证 \(U=\lambda^2/(2C')\)。
(3) 求介质内部束缚体电荷密度和两边界的束缚面电荷，并核验总束缚电荷守恒。
(4) 讨论 \(\alpha\to0\) 的均匀介质极限、介质极化的符号以及外导体电压变化方向。
